
\section{\BookRecordIndex{T0062}边界马斯洛夫指标（Joel Robbin 撰）}\label{sec:C-3}

本节定义定理~\ref{thm:C-1-10} 的指标公式中出现的\BookRecordIndex{T0062}边界马斯洛夫指标，
并建立其性质。基本方法是对曲面 $\Sigma$ 的裤形分解作归纳。
这类分解的存在是 Morse 理论的一个简单应用；借此可将相关结论归结为圆盘上的丛。

\paragraph{裤形分解。}
这里的二维流形是指可能带边界的紧致、定向、连通二维流形 $\Sigma$。

\begin{definition}\label{def:C-3-1}
二维流形 $\Sigma_{02}$ 的一个\emph{分解}，是它的一对子流形 $\Sigma_{01},\Sigma_{12}$，满足
\[
  \Sigma_{02}=\Sigma_{01}\cup\Sigma_{12},\qquad
  \Sigma_{01}\cap\Sigma_{12}=\partial\Sigma_{01}\cap\partial\Sigma_{12}.
\]
于是
\[
  \partial\Sigma_{ij}=\Gamma_i\cup\Gamma_j,\qquad
  \Gamma_i\cap\Gamma_j=\varnothing,
\]
其中 $\Gamma_i$ 是 $\Sigma_{02}$ 中有限个互不相交圆周的并，且
$\Gamma_1=\Sigma_{01}\cap\Sigma_{12}$。
\end{definition}

\begin{definition}\label{def:C-3-2}
\emph{带 $h$ 个洞的圆盘}，是从一个闭圆盘中去掉 $h$ 个互不相交闭子圆盘的内部
所得的带边二维流形。\emph{裤形曲面}是带两个洞的圆盘。
连通二维流形 $\Sigma$ 的一个\emph{裤形分解}，是子流形序列
\[
  \Sigma_0\subset\Sigma_1\subset\Sigma_2\subset\cdots\subset\Sigma_n=\Sigma,
\]
其中 $\Sigma_0$ 是圆盘，且对 $k=1,2,\ldots,n$，
$\Sigma_k$ 按定义~\ref{def:C-3-1} 分解为 $\Sigma_{k-1}$ 与一个裤形曲面
（在 $k=n$ 时也可能是一个圆盘）。
\end{definition}

可借 Morse 理论构造裤形分解。在闭曲面情形，存在一个 Morse 函数
$f:\Sigma\to[0,1]$，恰有一个局部极小点和一个局部极大点。
扰动 $f$，使每个临界水平集只含一个\BookRecordIndex{T0233}临界点，再沿一列正则水平集切开 $\Sigma$，
便得到一个始于圆盘、终于圆盘（分别是极小点和极大点的邻域）的分解，
每越过一个临界水平集就添加一个裤形曲面。若 $\Sigma$ 的边界非空，
选取 $f$ 使 $f^{-1}(1)=\partial\Sigma$，且 $1$ 是\BookRecordIndex{T0409}正则值，再照此进行。

\paragraph{裤形归纳。}
假设我们想证明关于二维流形、或给定二维流形的所有二维子流形的一个定理，
并已证明如下两步：
\begin{quote}
\textbf{基础步：}该定理对圆盘成立；

\textbf{归纳步：}若 $\Sigma_{02}$ 分解为 $\Sigma_{01}$ 与 $\Sigma_{12}$，
且该定理对三个流形 $\Sigma_{ij}$ 中的两个成立，则它对第三个也成立。
\end{quote}
那么该定理普遍成立。事实上，先按洞数归纳可得它对带洞圆盘成立；
再按裤形分解的长度 $n$ 归纳，可得它对至少有两个边界分支的流形成立。
最后，带边界的流形可由边界分支少一个的流形去掉一个圆盘得到，
因此该定理对所有流形成立。

\paragraph{马斯洛夫指标。}
我们回顾 Arnold 对马斯洛夫指标的定义（另一种处理见文献~[277] 第 2 章）。令
\[
  R(n):=\operatorname{GL}(n,\mathbb C)/\operatorname{GL}(n,\mathbb R)
\]
为 $\mathbb C^n$ 中全实子空间的流形。注意，$\mathbb R^{2n}=\mathbb C^n$
中拉格朗日子空间组成的流形
\sourcefootnote{原书把拉格朗日 Grassmann 流形还写成
$\operatorname{Sp}(2n)/\operatorname{GL}(n,\mathbb R)$。
这按通常的稳定子嵌入并非同一流形：当 $n=1$ 时该商为二维，
而 $L(1)=\operatorname U(1)/\operatorname O(1)$ 为一维。
辛群作用的完整稳定子还含对称矩阵的幂零部分；此处保留原式以待复核。}
\[
  L(n):=\operatorname U(n)/\operatorname O(n)
  =\operatorname{Sp}(2n)/\operatorname{GL}(n,\mathbb R)\subset R(n)
\]
是 $R(n)$ 的形变收缩核。定义 $\rho:R(n)\to S^1$ 为
\[
  \rho\bigl(a\cdot\operatorname{GL}(n,\mathbb R)\bigr)
  :=\frac{\det(a^2)}{\det(a^*a)},
\]
其中 $a^*$ 表示共轭转置。设 $\Gamma$ 是无边界的紧致定向一维流形，
即若干圆周的不交并。映射 $\Lambda:\Gamma\to R(n)$ 的\emph{马斯洛夫指标}定义为
\begin{equation}\label{eq:C-3-1}
  \mu(\Lambda):=\deg(\rho\circ\Lambda).
\end{equation}
当 $\Gamma$ 连通时，有如下结论。

\begin{theorem}[Arnold~\textup{[17]}]\label{thm:C-3-3}
$R(n)$ 中的两个圈同伦，当且仅当它们的马斯洛夫指标相同。
\end{theorem}

圆周 $S^1$ 上的复向量丛 $E$ 可被平凡化，因为 $\operatorname{GL}(n,\mathbb C)$ 连通。
当 $E=S^1\times\mathbb C^n$ 时，全实子丛 $F\subset E$ 通过
$\Lambda(z):=F_z$ 给出映射 $\Lambda:S^1\to R(n)$。
应区分 $\Lambda$ 与实向量丛 $F$。由于 $\operatorname{GL}(n,\mathbb R)$ 有两个连通分支，
圆周上有两类实秩 $n$ 向量丛。因此，当 $F_1,F_2$ 的马斯洛夫指标同奇偶时，
它们作为实向量丛同构。由 $E=F_j\oplus iF_j$，实同构 $F_1\to F_2$
可延拓为复同构；换言之，存在圈 $U:S^1\to\operatorname{GL}(n,\mathbb C)$ 使
\[
  \Lambda_2(z)=U(z)\Lambda_1(z).
\]
这样的圈 $U$ 可缩，当且仅当 $\mu(\Lambda_1)=\mu(\Lambda_2)$。

\begin{definition}\label{def:C-3-4}
$\Sigma$ 上的\emph{\BookRecordIndex{T0080}丛对} $(E,F)$，由复向量丛 $E\to\Sigma$ 和边界上的全实子丛
$F\subset E|_{\partial\Sigma}$ 构成。
$\Sigma_{02}$ 上\BookRecordIndex{T0080}丛对 $(E_{02},F_{02})$ 的一个\emph{分解}，由两个\BookRecordIndex{T0080}丛对
$(E_{01},F_0\cup F_1)$ 和 $(E_{12},F_1\cup F_2)$ 构成；前者在 $\Sigma_{01}$ 上，
后者在 $\Sigma_{12}$ 上，且 $\Sigma_{01},\Sigma_{12}$ 如定义~\ref{def:C-3-1}
分解 $\Sigma_{02}$，并有 $F_i\subset E_{02}|_{\Gamma_i}$。
\end{definition}

下一个定理公理化地定义\BookRecordIndex{T0062}边界马斯洛夫指标。

\begin{theorem}\label{thm:C-3-5}
存在唯一一种运算，称为\emph{\BookRecordIndex{T0062}边界马斯洛夫指标}，它给每个\BookRecordIndex{T0080}丛对 $(E,F)$
指定一个整数 $\mu(E,F)$，并满足以下公理：

\noindent\textup{（同构）} 若 $\Phi:E_1\to E_2$ 是覆盖保向微分同胚
$\phi:\Sigma_1\to\Sigma_2$ 的向量丛同构，则
\[
  \mu(E_1,F_1)=\mu(E_2,\Phi(F_1)).
\]
\noindent\textup{（直和）}
\[
  \mu(E_1\oplus E_2,F_1\oplus F_2)=\mu(E_1,F_1)+\mu(E_2,F_2).
\]
\noindent\textup{（复合）} 对定义~\ref{def:C-3-4} 的分解，有
\[
  \mu(E_{02},F_{02})=\mu(E_{01},F_{01})+\mu(E_{12},F_{12}).
\]
\noindent\textup{（归一化）} 当 $\Sigma=\mathbb D$ 是单位圆盘，
$E=\mathbb D\times\mathbb C$ 为平凡丛，且对
$z=e^{i\theta}\in\partial\mathbb D=S^1$ 有
$F_z=\mathbb R e^{ik\theta/2}$ 时，
\[
  \mu(\mathbb D\times\mathbb C,F)=k.
\]
\end{theorem}

下面列出其主要性质。回忆：在边界非空的连通黎曼曲面上，
每个复向量丛都可平凡化；证明见下面的引理~\ref{lem:C-3-8}。

\begin{theorem}\label{thm:C-3-6}
\BookRecordIndex{T0062}边界马斯洛夫指标满足：

\noindent\textup{（平凡丛）} 若 $\partial\Sigma\ne\varnothing$ 且
$E=\Sigma\times\mathbb C^n$，则
\begin{equation}\label{eq:C-3-2}
  \mu(\Sigma\times\mathbb C^n,F)=\mu(\Lambda),
\end{equation}
其中右侧由式~\eqref{eq:C-3-1} 定义，而 $\Lambda$ 定义为
\begin{equation}\label{eq:C-3-3}
  \Lambda(z):=F_z,\qquad z\in\partial\Sigma.
\end{equation}
\noindent\textup{（陈类）} 若 $\partial\Sigma=\varnothing$，
则 $\mu(E,\varnothing)$ 是第一陈类 $c_1(E)\in H^2(\Sigma)$
在基本类 $[\Sigma]\in H_2(\Sigma)$ 上取值的两倍：
\begin{equation}\label{eq:C-3-4}
  \mu(E,\varnothing)=2\langle c_1(E),[\Sigma]\rangle.
\end{equation}
\end{theorem}

\begin{theorem}\label{thm:C-3-7}
同一流形 $\Sigma$ 上的两个\BookRecordIndex{T0080}丛对 $(E_1,F_1)$ 与 $(E_2,F_2)$
在底空间恒等映射之上同构，当且仅当 $E_1,E_2$ 的秩相同，
$\mu(E_1,F_1)=\mu(E_2,F_2)$，且对 $\partial\Sigma$ 的每个分支 $C$，
实向量丛 $F_1|_C$ 与 $F_2|_C$ 同构。
\end{theorem}

定理~\ref{thm:C-3-5}、\ref{thm:C-3-6}、\ref{thm:C-3-7} 的证明要用到以下辅助引理。
设 $F\to\Gamma$ 是紧致一维流形上的实秩 $n$ 向量丛。
$F$ 的一个\emph{部分\BookRecordIndex{T0067}标架}是一个分解 $F=L\oplus F'$，
其中 $L\to\Gamma$ 为实线丛，$F'\to\Gamma$ 为子丛，并给定
$F'$ 的\BookRecordIndex{T0067}标架，即处处构成一组基的 $(n-1)$ 个截面
$s_2,\ldots,s_n:\Gamma\to F'$。

\begin{lemma}\label{lem:C-3-8}
设 $\partial\Sigma\ne\varnothing$，且 $\Sigma$ 上\BookRecordIndex{T0080}丛对 $(E,F)$
带有部分\BookRecordIndex{T0067}标架 $(L,F',s_2,\ldots,s_n)$。则存在平凡化
$\Phi:E\to\Sigma\times\mathbb C^n$ 和映射
$\lambda:\partial\Sigma\to S^1$，使得对 $z\in\partial\Sigma$，
\[
  \Phi(L_z)=\sqrt{\lambda(z)}\,\mathbb R e_1,\qquad
  \Phi(s_i(z))=e_i,\quad i=2,\ldots,n,
\]
其中 $e_1,\ldots,e_n$ 是 $\mathbb C^n$ 的标准基。特别地，
\[
  \Phi(F_z)=\sqrt{\lambda(z)}\,\mathbb R\times\mathbb R^{n-1},
  \qquad z\in\partial\Sigma.
\]
\end{lemma}

\begin{proof}
因 $\partial\Sigma\ne\varnothing$，$\Sigma$ 与若干圆周的楔和同伦等价。
又因 $\pi_0(\operatorname{GL}(n,\mathbb C))=0$，
任意复向量丛 $E\to\Sigma$ 都平凡，故可设 $E=\Sigma\times\mathbb C^n$。
按式~\eqref{eq:C-3-3} 定义 $\Lambda$，并选取 $\lambda$，使对
$\partial\Sigma$ 的每个分支 $C$，
$\mu(\Lambda|_C)=\mu(\Lambda'|_C)$，其中
$\Lambda'(z)=\sqrt{\lambda(z)}\,\mathbb R\times\mathbb R^{n-1}$。
由 Arnold 定理，$\Lambda$ 与 $\Lambda'$ 同伦。取光滑映射
$U:\Gamma\to\operatorname{GL}(n,\mathbb C)$，使得对每个 $z\in\partial\Sigma$，
\[
  U(z)L_z=\sqrt{\lambda(z)}\,\mathbb R e_1,\qquad
  U(z)s_i(z)=e_i,\quad i=2,\ldots,n.
\]
则 $U(z)\Lambda(z)=\Lambda'(z)$。再由 Arnold 定理，
$U$ 在 $\partial\Sigma$ 每个分支上的限制可缩。
故 $U$ 可延拓为 $\partial\Sigma$ 的领状邻域上的映射（仍记 $U$），
且在该领状邻域的内侧边界上为恒等映射。再在 $\Sigma$ 其余部分将 $U$
延拓为恒等映射，并定义 $\Phi(z,v)=(z,U(z)v)$。
\end{proof}

\begin{corollary}\label{cor:C-3-9}
圆盘 $\mathbb D$ 上的任意秩一\BookRecordIndex{T0080}丛对 $(E,F)$ 都有平凡化
$\Phi:E\to\mathbb D\times\mathbb C$，使得对
$z=e^{i\theta}\in\partial\mathbb D=S^1$，
$\Phi(F_z)=\mathbb R e^{ik\theta/2}$。
\end{corollary}
\begin{proof}
在引理~\ref{lem:C-3-8} 的证明中取 $\lambda(e^{i\theta})=e^{ik\theta}$ 即得。
\end{proof}

\begin{proof}[定理~\ref{thm:C-3-5} 与定理~\ref{thm:C-3-6} 的证明]
定理~\ref{thm:C-3-5} 的唯一性直接由各公理、
推论~\ref{cor:C-3-9} 和裤形归纳\BookRecordIndex{T0326}推出。为证存在性，
按定理~\ref{thm:C-3-6} 所要求的方式定义 $\mu$，再验证各公理。
先须说明：若 $\partial\Sigma\ne\varnothing$ 且 $(E,F)$ 是 $\Sigma$ 上的\BookRecordIndex{T0080}丛对，则
\begin{equation}\label{eq:C-3-5}
  \mu(E,F)=\mu(\Lambda_\Phi),\qquad
  \Lambda_\Phi(z)=\Phi(F_z),
\end{equation}
不依赖所用平凡化 $\Phi:E\to\Sigma\times\mathbb C^n$。
另一平凡化 $\Psi:E\to\Sigma\times\mathbb C^n$ 确定
$U:\Sigma\to\operatorname{GL}(n,\mathbb C)$，满足
\[
  (\Phi\circ\Psi^{-1})(z,v)=(z,U(z)v),\quad z\in\Sigma,\ v\in\mathbb C^n,
\]
且 $U(z)\Lambda_\Psi(z)=\Lambda_\Phi(z)$，$z\in\partial\Sigma$。
若 $a\in\operatorname{GL}(n,\mathbb C)$ 满足
$a(\mathbb R^n)=\Lambda_\Psi(z)$，则
$(U(z)a)(\mathbb R^n)=\Lambda_\Phi(z)$。
故由式~\eqref{eq:C-3-1}，
\begin{equation}\label{eq:C-3-6}
  \rho(\Lambda_\Phi(z))=g(z)\rho(\Lambda_\Psi(z)),\qquad
  g(z)=\frac{\det(U(z)^2)}{\det(U(z)^*U(z))},
\end{equation}
$z\in\partial\Sigma$。但 $g:\partial\Sigma\to S^1$ 可延拓到 $\Sigma$，
故其次数为零。因此
$\mu(\Lambda_\Phi)=\deg(\rho\circ\Lambda_\Phi)
=\deg(\rho\circ\Lambda_\Psi)=\mu(\Lambda_\Psi)$，如所断言。

现可同时证明定理~\ref{thm:C-3-5} 的存在性和定理~\ref{thm:C-3-6}：
当 $\partial\Sigma\ne\varnothing$ 时按式~\eqref{eq:C-3-2}，
当 $\partial\Sigma=\varnothing$ 时按式~\eqref{eq:C-3-4} 定义
$\mu(E,F)$，并验证它满足定理~\ref{thm:C-3-5} 的公理。
除分解公理外，各公理都直接成立。计算 $\rho\circ\Lambda$ 的次数时，
默认 $\partial\Sigma$ 取从 $\Sigma$ 诱导的定向。因此对于分解公理中的分解，有
\[
  \mu(\Lambda_{02})=\mu(\Lambda_{01})+\mu(\Lambda_{12}),\qquad
  \Lambda_{ij}(z)=(F_{ij})_z,
\]
因为交界 $\Sigma_{01}\cap\Sigma_{12}$ 上的贡献相互抵消。
当 $\partial\Sigma_{02}\ne\varnothing$ 时，这就证明了分解公理。
当 $\partial\Sigma_{02}=\varnothing$ 时还须进一步论证；
为简化记号，将其改写为下面的定理。
\end{proof}

\begin{theorem}\label{thm:C-3-10}
设 $\partial\Sigma=\varnothing$，且 $\Sigma=\Sigma_0\cup\Sigma_1$
是定义~\ref{def:C-3-1} 意义下的分解，记
$\Gamma:=\Sigma_0\cap\Sigma_1=\partial\Sigma_0=\partial\Sigma_1$。
设 $E$ 是 $\Sigma$ 上的向量丛，$F\subset E|_\Gamma$ 是全实子丛，
并记 $E_i=E|_{\Sigma_i}$。则
\[
  2\langle c_1(E),[\Sigma]\rangle
  =\mu(E_0,F)+\mu(E_1,F).
\]
\end{theorem}

\begin{proof}
由引理~\ref{lem:C-3-8}，$E$ 分裂为一个线丛与平凡丛
$\Sigma\times\mathbb C^{n-1}$，且
$F\cap(\Gamma\times\mathbb C^{n-1})=\Gamma\times\mathbb R^{n-1}$。
故不失一般性可设 $n=1$。再设 $E$ 带有埃尔米特形式。
\BookRecordIndex{T0391}酉平凡化 $\Phi_i:E|_{\Sigma_i}\to\Sigma_i\times\mathbb C$
可写成 $\Phi_i^{-1}(z,v)=v s_i(z)$，其中 $s_i:\Sigma_i\to E_i$ 是范数为一的截面。
定义 $u:\Gamma\to S^1$ 为
\[
  s_1(z)=u(z)s_0(z),\qquad z\in\Gamma.
\]
并令 $\Lambda_i(z):=\Phi_i(F_z)\subset\mathbb C$。
于是 $\Lambda_0(z)=u(z)\Lambda_1(z)$，故如式~\eqref{eq:C-3-6} 那样，
$\rho(\Lambda_0)=u^2\rho(\Lambda_1)$。
现以 $\Sigma_0$ 的边界定向来定向 $\Gamma$，则
$\mu(E_0,F)=\deg(\rho\circ\Lambda_0)$，
$\mu(E_1,F)=-\deg(\rho\circ\Lambda_1)$。
因此，取次数得
\begin{equation}\label{eq:C-3-7}
  \mu(E_0,F)+\mu(E_1,F)=2\deg(u).
\end{equation}
取 $E$ 上\BookRecordIndex{T0052}埃尔米特联络 $\nabla$，并定义
\[
  \alpha_0:=s_0^{-1}\nabla s_0\in\Omega^1(\Sigma_0,i\mathbb R),\qquad
  \alpha_1:=s_1^{-1}\nabla s_1\in\Omega^1(\Sigma_1,i\mathbb R).
\]
则 $F^\nabla|_{\Sigma_0}=d\alpha_0$、
$F^\nabla|_{\Sigma_1}=d\alpha_1$，且
$(\alpha_1-\alpha_0)|_\Gamma=u^{-1}du$。故
\begin{align*}
  \langle c_1(E),[\Sigma]\rangle
  &=\frac{i}{2\pi}\int_\Sigma F^\nabla\\
  &=\frac{i}{2\pi}\int_{\Sigma_0}d\alpha_0
    +\frac{i}{2\pi}\int_{\Sigma_1}d\alpha_1\\
  &=\frac{i}{2\pi}\int_\Gamma(\alpha_0-\alpha_1)\\
  &=\frac{1}{2\pi i}\int_\Gamma u^{-1}du
  =\deg(u).
\end{align*}
由式~\eqref{eq:C-3-7} 即得结论。
\end{proof}

\begin{proof}[定理~\ref{thm:C-3-7} 的证明]
如定理~\ref{thm:C-3-10} 的证明，不失一般性可设秩为一。
先设 $\partial\Sigma\ne\varnothing$，并设
$E_1=E_2=\Sigma\times\mathbb C$。由假设可取
$u:\partial\Sigma\to\operatorname{GL}(1,\mathbb C)$，使得
$u(z)F_1(z)=F_2(z)$；须证明 $u$ 可延拓到 $\Sigma$。
不失一般性设 $|u|\equiv1$。由假设，$u$ 的总次数为零，
但它在单个边界分支上的限制未必次数为零。按边界分支数归纳：
若只有一个分支，则 $u$ 同伦于常值映射，利用该同伦可在领状邻域延拓，
再把其内侧边界上的常值映射延拓到 $\Sigma$ 其余部分。
若某个分支 $C$ 上 $u|_C$ 的次数为零，则可在该分支填入圆盘，
并将 $u$ 延拓到该圆盘，从而减少边界分支数。
若每个边界分支 $C$ 上 $u|_C$ 的次数均非零，
则不同的两个边界分支上存在 $u$ 取相同值的两点。
用一条弧连接这两点，在该弧上将 $u$ 延拓为公共值，
再沿弧切开，便减少了边界分支数。

$\partial\Sigma=\varnothing$ 的情形是熟知的：线丛由第一陈类刻画。
也可将 $\Sigma$ 分解为 $\Sigma'$ 与圆盘 $\mathbb D$，
引入全实子丛 $F_i\subset E_i|_{\partial\mathbb D}$，使
$\mu(E_1|_{\mathbb D},F_1)=\mu(E_2|_{\mathbb D},F_2)$，
再将边界非空的情形应用于 $E_i|_{\Sigma'}$。
利用定理~\ref{thm:C-3-10} 的证明，将 $c_1(E_i)$
与 $\mu(E_i|_{\mathbb D},F_i)$ 联系起来。
\end{proof}
